\documentclass[../main.tex]{subfiles}
\begin{document}

\chapter{コンパイラによる ILP}
\lessoninfo{
  video={Lesson 10，動画 1--20},
  textbook={H\&P \cite{hennessy2017} §3.2，付録 H},
  keywords={命令レベル並列性，サイクル当たり命令数，依存の連鎖，木の高さ削減，命令スケジューリング，インオーダ実行，if 変換，ループ展開，インライン展開，コード肥大化，命令数，ソフトウェアパイプライニング，トレーススケジューリング},
}

第6〜9章では，ハードウェアによってプロセッサのサイクル当たり命令数（IPC）
を高める手法を見てきた．レジスタリネーミング，リオーダバッファを伴う
アウトオブオーダ実行，ロードとストアのアウトオブオーダ実行により，IPC は
プログラムの命令レベル並列性（ILP）に近づく．しかしハードウェアが利用できる
並列性は，プログラムが実際に含んでいて，かつハードウェアの手の届く範囲に
あるものに限られる．この章では，コンパイラがその両方をどのように改善するかを
示す．木の高さ削減は ILP を制限する依存の連鎖を短くし，命令スケジューリングは
独立な命令を隣り合わせに配置し，if 変換・ループ展開・インライン展開は，
それらを隔てる分岐・呼び出し・リターンを除去する．\caution{「コンパイラ ILP」は
2つ目の ILP ではない．ILP はプログラムの性質であり（第6章），コンパイラはその
プログラムを変えるのである．}

\section{コンパイラはどのように IPC を高められるか}
\label{sec:l10-limits}

どのプロセッサでも，IPC は実行するプログラムの ILP によって上限が定まる
（\cref{sec:l06-ipc}）．\source{Lesson 10，動画 1--2；H\&P §3.2．}IPC に
対する 2つの制限はプロセッサだけで決まるものではなく，コンパイラはその
どちらにも影響する．
\begin{description}
  \item[プログラムの ILP] ILP は，命令数を RAW 依存の最長の連鎖の長さで
    割ったものである（\cref{eq:l06-ilp,eq:l06-cycle}）．連鎖がどれだけ長く
    なるかは，コンパイラがソースプログラムをどう翻訳したかに依存する．同じ
    計算を，より短い連鎖をなす命令で表現できることが多い．
  \item[プロセッサが見つけられる並列性] プロセッサが実行する命令を探すのは，
    最も古い未実行の命令に続く限られた数の命令の中だけである．アウトオブ
    オーダプロセッサはリオーダバッファに収まる命令を対象とし（第8章），
    インオーダプロセッサは次の数命令だけを見て，待たなければならない最初の
    命令で止まる（第6章）．プログラム順でこのウィンドウより離れた位置に
    ある独立な命令は，ILP がどれほど高くても同じサイクルには実行できない．
    コンパイラはそのような命令を近くに配置でき，また両者の間にある分岐，
    呼び出し，リターンを除去できる．
\end{description}
コンパイラはプログラムの一部のウィンドウではなく全体を見ることができ，
しかもプログラムの実行前に動作するので，プロセッサが 1 クロックサイクルの
中では到底行えない解析にも時間をかけられる．\margin{コンパイラのウィンドウは
関数全体であり，インライン展開後は呼び出し木全体である．プロセッサのウィンドウは
リオーダバッファである（第8章）．}この章で扱う技法を
\cref{tab:l10-overview} に示す．

\begin{table}[htbp]
  \centering\small
  \caption{IPC に対する制限と，それに対処するコンパイラ技法．}
  \label{tab:l10-overview}
  \begin{tabular}{@{}l>{\raggedright\arraybackslash}p{0.35\linewidth}c@{}}
    \toprule
    制限 & コンパイラ技法 & 節 \\
    \midrule
    長い依存の連鎖が ILP を下げる      & 木の高さ削減
      & \ref{sec:l10-thr} \\
    独立な命令が離れている             & 命令スケジューリング
      & \ref{sec:l10-scheduling} \\
    分岐が命令を隔てる                 & if 変換，\newline トレーススケジューリング
      & \ref{sec:l10-if-conversion}，\ref{sec:l10-other} \\
    ループの分岐が反復を隔てる         & ループ展開，\newline ソフトウェアパイプライニング
      & \ref{sec:l10-unrolling}，\ref{sec:l10-other} \\
    呼び出しとリターンが命令を隔てる   & インライン展開
      & \ref{sec:l10-inlining} \\
    \bottomrule
  \end{tabular}
\end{table}

\section{木の高さ削減}
\label{sec:l10-thr}

コンパイラは算術式を\emph{式木}を経由して翻訳する．式木では，各演算が
ノードであり，その子がオペランドである．\source{Lesson 10，動画 3--4；
H\&P 付録 H．}各演算は 1つの命令になり，あるノードの命令はその子の命令の
結果を読む．したがって，これらの命令の間の RAW 依存の最長の連鎖は木の
\emph{高さ}，すなわち根から葉への最長経路上の演算の個数に等しく，第6章の
理想プロセッサでは木の 1 段が 1 サイクルを要する．\margin{1 段 1 サイクルは，
すべての演算が 1 サイクルであることを仮定している．多サイクルの演算では 1 段が
そのレイテンシ分かかるので，高さの影響はさらに大きい．}

素直な翻訳は，ソースプログラムがオペランドを書いた順に従う．R1 から R6 に
保持された 6つの変数の和を R10 に求める場合，次のコードが生成される．
\begin{lstlisting}[language={[x86masm]Assembler}, numbers=none]
ADD R10, R1, R2      ; サイクル 1
ADD R10, R10, R3     ; サイクル 2
ADD R10, R10, R4     ; サイクル 3
ADD R10, R10, R5     ; サイクル 4
ADD R10, R10, R6     ; サイクル 5
\end{lstlisting}
どの加算も直前の加算の結果を使う．木は高さ~5 の連鎖に退化し
（\cref{fig:l10-tree-height}a），5つの命令はどのプロセッサでも 5 サイクル
かかり，その ILP は~1 である．

\begin{definition}[木の高さ削減]
  \label{def:l10-thr}
  演算の結合法則と交換法則を用いて，式を，より段数の少ない式木を持つ等価な
  式に書き換えることを\term[きのたかささくげん]{木の高さ削減}[tree height reduction]%
  という．演算の個数，したがって命令数は変わらない．
\end{definition}

\begin{example}
  \label{ex:l10-thr}
  $((\text{R1}+\text{R2}) + (\text{R3}+\text{R4})) +
  (\text{R5}+\text{R6})$ とグループ化すると，同じ和は次のようになる．
  \begin{lstlisting}[language={[x86masm]Assembler}, numbers=none]
ADD R10, R1, R2      ; サイクル 1
ADD R11, R3, R4      ; サイクル 1
ADD R12, R5, R6      ; サイクル 1
ADD R10, R10, R11    ; サイクル 2
ADD R10, R10, R12    ; サイクル 3
\end{lstlisting}
  対をなす 3つの加算はどれも何にも依存せず，サイクル~1 で実行される．
  最初の 2 対の和がサイクル~2 に，最後の加算がサイクル~3 に続く
  （\cref{fig:l10-tree-height}b）．
  命令は依然として 5つだが，それが 3 サイクルで済む．ILP は 1 から
  $5/3 \approx 1.67$ に上がる．その代償は，中間結果のための 2つの追加の
  レジスタ R11 と R12 である．\margin{再グループ化の予算はレジスタである．
  アーキテクチャ整数レジスタは x86-64 が 16 本，A64 と RISC-V が 31 本であり，
  それを超えるとコンパイラはスタックへ退避する．}
\end{example}

\begin{figure}[htbp]
  \centering
  % Tree height reduction of a six-operand sum: (a) left-to-right translation,
% a chain of height 5; (b) balanced tree of height 3.  Each level of
% operations is executed in one cycle of the ideal processor.
\begin{tikzpicture}[x=1cm,y=1cm, font=\footnotesize\sffamily,
  op/.style={circle, draw, fill=ln-box, minimum size=0.5cm, inner sep=0pt},
  leaf/.style={font=\scriptsize\ttfamily, inner sep=1pt, fill=white},
  edge/.style={thick, ln-accent},
  res/.style={font=\scriptsize\ttfamily, inner sep=1pt}]
  % ---- cycle rows
  \foreach \c in {1,...,5} {
    \draw[ln-rule, dotted] (-0.2,{(\c-1)*0.8}) -- (9.7,{(\c-1)*0.8});
    \node[anchor=east, text=ln-muted, font=\scriptsize\sffamily]
      at (-0.25,{(\c-1)*0.8}) {\iflangja{サイクル \c}{cycle \c}};
  }
  % ---- (a) left to right: ((((R1+R2)+R3)+R4)+R5)+R6
  \node[leaf] (r1) at (0.2,-0.8) {R1};
  \node[leaf] (r2) at (1.0,-0.8) {R2};
  \node[op] (a1) at (0.6,0.0) {$+$};
  \node[op] (a2) at (1.2,0.8) {$+$};
  \node[op] (a3) at (1.8,1.6) {$+$};
  \node[op] (a4) at (2.4,2.4) {$+$};
  \node[op] (a5) at (3.0,3.2) {$+$};
  \node[leaf] (r3) at (1.8,0.0) {R3};
  \node[leaf] (r4) at (2.4,0.8) {R4};
  \node[leaf] (r5) at (3.0,1.6) {R5};
  \node[leaf] (r6) at (3.6,2.4) {R6};
  \draw[edge] (a1) -- (r1) (a1) -- (r2)
              (a2) -- (a1) (a2) -- (r3)
              (a3) -- (a2) (a3) -- (r4)
              (a4) -- (a3) (a4) -- (r5)
              (a5) -- (a4) (a5) -- (r6);
  \node[res] (ra) at (3.0,3.85) {R10};
  \draw[edge] (a5) -- (ra);
  \node at (1.9,-1.45) {\iflangja{(a) 左から順に：高さ 5}{(a) left to right: height 5}};
  % ---- (b) balanced: ((R1+R2)+(R3+R4))+(R5+R6)
  \begin{scope}[xshift=5.6cm]
    \node[leaf] (s1) at (0.0,-0.8) {R1};
    \node[leaf] (s2) at (0.9,-0.8) {R2};
    \node[leaf] (s3) at (1.5,-0.8) {R3};
    \node[leaf] (s4) at (2.4,-0.8) {R4};
    \node[leaf] (s5) at (3.0,-0.8) {R5};
    \node[leaf] (s6) at (3.9,-0.8) {R6};
    \node[op] (p1) at (0.45,0.0) {$+$};
    \node[op] (p2) at (1.95,0.0) {$+$};
    \node[op] (p3) at (3.45,0.0) {$+$};
    \node[op] (p4) at (1.2,0.8) {$+$};
    \node[op] (p5) at (2.3,1.6) {$+$};
    \draw[edge] (p1) -- (s1) (p1) -- (s2)
                (p2) -- (s3) (p2) -- (s4)
                (p3) -- (s5) (p3) -- (s6)
                (p4) -- (p1) (p4) -- (p2)
                (p5) -- (p4) (p5) -- (p3);
    \node[res] (rb) at (2.3,2.25) {R10};
    \draw[edge] (p5) -- (rb);
    \node at (1.95,-1.45) {\iflangja{(b) 平衡木：高さ 3}{(b) balanced: height 3}};
  \end{scope}
\end{tikzpicture}

  \caption{R1 から R6 までの和の式木．(a) 左から順に翻訳すると，5つの加算の
    連鎖になる．(b) \cref{ex:l10-thr} の平衡木は高さ~3 である．加算の各段は
    理想プロセッサの 1 サイクルで実行される．}
  \label{fig:l10-tree-height}
\end{figure}

1つの結合的な演算で結ばれた $n$ 個のオペランドについて，左から順に翻訳
すると高さは $n-1$ になるが，平衡木の高さは $\lceil \log_2 n \rceil$ で
ある．加算と乗算は結合的であり，AND，OR，XOR も同様である．
\caution{浮動小数点の加算は（丸めのため）結合的ではない．コンパイラがこれを
再グループ化するのは，\texttt{-ffast-math} のように許可された場合だけである．}
減算は $(a-b)-c \ne a-(b-c)$ なので結合的ではないが，符号を調整すれば
減算を含む式も再グループ化できる．$a - b = a + (-b)$ であるから，各項が
自身の符号を保つ限り項は自由にグループ化でき，マイナス符号に続くグループの
中の符号は反転する．例えば $a-b-c-d = (a-b)-(c+d)$ となり，高さが 3 から~2
に減る．

\begin{keypoint}
  木の高さ削減は，命令間の依存を変えるのであって命令数を変えるのではない．
  ILP を制限する連鎖を短くするので，1 サイクルに複数の命令を実行できる
  プロセッサで効果がある．
\end{keypoint}

\section{命令スケジューリング}
\label{sec:l10-scheduling}

プロセッサが独立な命令を同じサイクルに実行するのは，それらを見つけられた
場合だけである．\source{Lesson 10，動画 5--7；H\&P §3.2．}コンパイラは通常
プログラムを文ごとに翻訳するので，ある文の命令の後に次の文の命令が続く．
各文が 1つの連鎖をなす場合，異なる文に属する独立な命令は連鎖の残りによって
隔てられる．ウィンドウが両方の連鎖を覆うアウトオブオーダプロセッサならそれ
でも見つけられるが，ウィンドウがより短いアウトオブオーダプロセッサは見つけ
られず，待たなければならない最初の命令までしか後続を実行しないインオーダ
プロセッサも同様である．\caution{メモリもコンパイラを縛る．ロードをストアより
前に移動できるのは，2つのアドレスが異なると証明できる場合だけである．同じ型の
ポインタ 2つでは証明できず，順序はそのままになる．}独立な命令を見つけやすくする
コンパイラ技法が 3つ
ある．命令スケジューリング，ループ展開（\cref{sec:l10-unrolling}），
トレーススケジューリング（\cref{sec:l10-other}）である．

\begin{definition}[命令スケジューリング]
  \label{def:l10-scheduling}
  プログラムの結果を変えずに命令を並べ替え，同じサイクルに実行できる命令が
  プログラム順で隣り合うようにすることを
  \term[めいれいすけじゅーりんぐ]{命令スケジューリング}[instruction scheduling]%
  という．
\end{definition}

インオーダプロセッサでは，スケジューリングは，そのままではプロセッサが待つ
ことになるサイクルを，待っている命令に依存しない命令で埋める．1 サイクルに
最大 $w$ 命令を実行するプロセッサ向けにスケジューリングするには，コンパイラ
は各命令について理想プロセッサがそれを実行するサイクルを求め
（\cref{eq:l06-cycle}），そのサイクルの順に，高々 $w$ 命令ずつのグループに
して命令を並べればよい．このとき各命令は，自分が読む結果を生成する命令より
後に置かれたままであるが，それに加えて，並べ替えたプログラムは元の
プログラムと同じ結果を計算しなければならない．WAW 依存や WAR 依存を破る
移動は，ある命令が使う値を変えてしまう．そこでコンパイラは，元の順序を
保つか，あるいは次のように補正する．
\begin{description}
  \item[宛先レジスタ] 移動する命令が，後の命令がまだ読む古い値を持つ
    レジスタを上書きしてしまう場合，その命令は代わりに，使われていない別の
    レジスタに結果を書き，その結果を読む命令もそれに合わせて変更する．
  \item[アドレスのオフセット] 添字の増分のようにレジスタに定数を加える命令
    が，そのレジスタをアドレスのベースとして使う命令より前に移動する場合，
    アドレスの変位をその定数だけ減らす．\texttt{SW R3, 0(R1)} は
    \texttt{ADDI R1, R1, 4} の後に置かれると \texttt{SW R3, -4(R1)} になる．%
    \tip{ある命令を，$c$ を加える \texttt{ADDI} より後ろへ移すと，そのレジスタを
    使う変位はすべて $c$ 減る．逆向きに移せば $c$ 増える．}
\end{description}

\begin{example}
  \label{ex:l10-scheduling}
  $a$ から $h$ を R1 から R8 に保持するとき，文 $x = (a+b)\times c - d$ と
  $y = (e-f)\times g + h$ は次のように翻訳される．
  \begin{lstlisting}[language={[x86masm]Assembler}, numbers=none]
I1:  ADD R9, R1, R2      ; a + b
I2:  MUL R9, R9, R3      ; (a + b) * c
I3:  SUB R9, R9, R4      ; x <- (a + b) * c - d
I4:  SUB R10, R5, R6     ; e - f
I5:  MUL R10, R10, R7    ; (e - f) * g
I6:  ADD R10, R10, R8    ; y <- (e - f) * g + h
\end{lstlisting}
  これらの命令は I1--I3 と I4--I6 の 2つの連鎖をなすので，理想プロセッサは
  サイクル~1 に I1 と I4 を，サイクル~2 に I2 と I5 を，サイクル~3 に I3 と
  I6 を実行する．ILP は~2 である．幅 2 のインオーダプロセッサは，I2 が I1
  を待たなければならないのでサイクル~1 には I1 だけを，サイクル~2 には I2
  だけを実行する．サイクル~3 では I3 を I4 と一緒に実行するが，その後は
  I5 と I6 が再び 1 サイクルずつ必要になる．合計 5 サイクル，IPC は
  $6/5 = 1.2$ である（\cref{fig:l10-scheduling}a）．I1，I4，I2，I5，I3，I6
  の順にスケジューリングすると，隣り合う命令の組はすべて独立になり，同じ
  プロセッサが 3 サイクルで済む．IPC は~2 であり，ILP に等しい
  （\cref{fig:l10-scheduling}b）．2つの文は異なるレジスタを使うので，この
  並べ替えは WAW 依存も WAR 依存も破らない．
\end{example}

\begin{figure}[htbp]
  \centering
  % Instruction scheduling for a two-wide in-order processor: (a) the two
% statements compiled one after the other; (b) the scheduled order.
% Arcs: RAW dependences (first statement above, second below); shaded
% groups: instructions executed in the same cycle; numbers: that cycle.
\begin{tikzpicture}[x=1cm,y=1cm, font=\footnotesize\sffamily, >={Stealth[length=1.6mm]},
  ins/.style={draw, fill=white, minimum width=0.9cm, minimum height=0.5cm, inner sep=0pt},
  dep/.style={->, thick, ln-accent},
  grp/.style={fill=ln-accent!12, rounded corners=2pt},
  cyc/.style={font=\scriptsize\sffamily, text=ln-muted}]
  % ---------------- (a) statement by statement
  \node[anchor=west] at (0.2,1.0)
    {\iflangja{(a) 文ごとに翻訳した順序：6 命令を 5 サイクルで実行}%
              {(a) statement by statement: 6 instructions in 5 cycles}};
  \foreach \xa/\xb in {1.8/1.8, 3.3/3.3, 4.8/6.3, 7.8/7.8, 9.3/9.3}
    \fill[grp] (\xa-0.6,-0.4) rectangle (\xb+0.6,0.4);
  \foreach \n/\x in {1/1.8, 2/3.3, 3/4.8, 4/6.3, 5/7.8, 6/9.3}
    \node[ins] (a\n) at (\x,0) {I\n};
  \draw[dep] (a1.north) to[bend left=40] (a2.north);
  \draw[dep] (a2.north) to[bend left=40] (a3.north);
  \draw[dep] (a4.south) to[bend right=40] (a5.south);
  \draw[dep] (a5.south) to[bend right=40] (a6.south);
  \node[cyc, anchor=east] at (1.1,-1.0) {\iflangja{サイクル}{cycle}};
  \foreach \c/\x in {1/1.8, 2/3.3, 3/4.8, 3/6.3, 4/7.8, 5/9.3}
    \node[cyc] at (\x,-1.0) {\c};
  % ---------------- (b) scheduled
  \begin{scope}[yshift=-2.8cm]
    \node[anchor=west] at (0.2,1.0)
      {\iflangja{(b) スケジューリング後：6 命令を 3 サイクルで実行}%
                {(b) scheduled: 6 instructions in 3 cycles}};
    \foreach \xa/\xb in {1.8/3.3, 4.8/6.3, 7.8/9.3}
      \fill[grp] (\xa-0.6,-0.4) rectangle (\xb+0.6,0.4);
    \foreach \n/\x in {1/1.8, 4/3.3, 2/4.8, 5/6.3, 3/7.8, 6/9.3}
      \node[ins] (b\n) at (\x,0) {I\n};
    \draw[dep] (b1.north) to[bend left=25] (b2.north);
    \draw[dep] (b2.north) to[bend left=25] (b3.north);
    \draw[dep] (b4.south) to[bend right=25] (b5.south);
    \draw[dep] (b5.south) to[bend right=25] (b6.south);
    \node[cyc, anchor=east] at (1.1,-1.0) {\iflangja{サイクル}{cycle}};
    \foreach \c/\x in {1/1.8, 1/3.3, 2/4.8, 2/6.3, 3/7.8, 3/9.3}
      \node[cyc] at (\x,-1.0) {\c};
  \end{scope}
\end{tikzpicture}

  \caption{幅 2 のインオーダプロセッサにおける \cref{ex:l10-scheduling} の
    プログラム．弧は RAW 依存であり，網掛けのグループは同じサイクルに実行
    される．(a) 文の順序のままでは，連鎖 I1--I3 が I4 を足止めする．
    (b) スケジューリング後は，両方の連鎖が毎サイクル進む．}
  \label{fig:l10-scheduling}
\end{figure}

\begin{example}
  \label{ex:l10-loop}
  次のループは
  \begin{lstlisting}[language=C, numbers=none]
for (i = 0; i < n; i++)
    a[i] = a[i] + x;
\end{lstlisting}
  整数配列の $n$ 個の要素それぞれに $x$ を加える．R1 が現在の要素を指し，
  R2 が $x$ を保持し，R5 が最後の要素の次のアドレスを保持するとき，次の
  ように翻訳される．\margin{翻訳後のループに $i$ はない．コンパイラは帰納変数
  としてアドレスだけを残し，それを R5 の終端と比較する（\emph{強度低減}）．}
  \begin{lstlisting}[language={[x86masm]Assembler}, numbers=none]
Loop: LW   R3, 0(R1)      ; R3 <- a[i]
      ADD  R3, R3, R2     ; R3 <- a[i] + x
      SW   R3, 0(R1)      ; a[i] <- R3
      ADDI R1, R1, 4      ; 次の要素へ
      BNE  R1, R5, Loop   ; 末尾まで繰り返す
\end{lstlisting}
  完全な分岐予測を持つ幅 4 のインオーダプロセッサは，ADD が LW を待つので
  まず LW だけを，次に SW が ADD を待つので ADD だけを実行する．SW と ADDI
  は一緒に実行され，ADDI を待つ BNE は次の反復の LW と一緒に実行される．
  1 反復は 3 サイクルかかり，CPI は $3/5 = 0.6$ である．ADDI を LW の直後に
  繰り上げると，ADDI はもはや連鎖の後ろで待たされず，増分の後に置かれること
  になった SW は変位 $-4$ を必要とする．
  \begin{lstlisting}[language={[x86masm]Assembler}, numbers=none]
Loop: LW   R3, 0(R1)
      ADDI R1, R1, 4
      ADD  R3, R3, R2
      SW   R3, -4(R1)     ; R1 はすでに進んでいる
      BNE  R1, R5, Loop
\end{lstlisting}
  これは理想サイクルの順序であり，分岐だけは反復の末尾に留めてある．%
  \caution{$-4$ は要素数ではなくバイト数である．変位は，要素の大きさによらず，
  ベースレジスタが増えた定数ちょうどで補正する．}LW と
  ADDI が一緒に実行され，ADD が続き，SW と BNE，そして次の反復の LW と ADDI
  が次のサイクルの 4つの枠をすべて埋める．1 反復は 2 サイクルとなり，CPI は
  $2/5 = 0.4$ である．
\end{example}

このようにスケジューリングが効くのは，\margin{静的なスケジュールは仮定した
レイテンシに固定される．ロードがキャッシュにヒットする限り順序は正しいが，ミス
するとインオーダプロセッサは止まり，反応できるのは第7〜9章のハードウェアだけで
ある．}インオーダプロセッサ，プログラム中の
連鎖に比べてウィンドウが短いアウトオブオーダプロセッサ，そして幅より多くの
命令が準備できているときのアウトオブオーダプロセッサである．どの命令が先に
実行されるかを順序が決めるからである（第6章）．

\section{スケジューリングと if 変換}
\label{sec:l10-if-conversion}

命令スケジューリングが命令を動かせるのは，全体がひとまとまりに実行される
コードの範囲の中だけであり，条件分岐はその範囲を終わらせる．
\source{Lesson 10，動画 8--9；H\&P §3.2．}分岐に続く命令は，その一方の経路
でしか実行されない．それを分岐の前に移動すれば，もう一方の経路でも実行され
てしまう．また 2つの経路の命令は，どちらか一方の経路しか実行されないので，
同じグループに置くことはできない．したがってコンパイラは，分岐の前後の命令
が独立であっても，分岐をまたいで命令をグループ化できない．\margin{この範囲を
\emph{基本ブロック}という．整数コードでは 5 命令に 1 回ほど分岐があり（第3章），
1つのブロックの中だけではスケジューリングの余地は小さい．}

小さな if-then-else に対しては，if 変換\index{if へんかん@if 変換}（第5章）
がこの境界を取り除く．変換後のコードは，条件と両方の部分を無条件に計算し，
条件付きムーブで結果を選択する．全体が分岐のない 1つの列になっており，
コンパイラはそれを他のコードと同じようにスケジューリングできる．条件と
2つの部分を計算する命令は選択に依存しないので，それら自身の依存に従って
先頭のグループに置かれ，選択がその後に続く．\cref{ex:l05-movzn} の if 変換
後のコードは，すでにこの形をしている．最初の 4つの命令は互いに独立であり，
3つの条件付きムーブはそれらにしか依存しないので，幅 4 のインオーダ
プロセッサは 7つの命令すべてを 2 サイクルで実行する．

\subsection{ループ}

ループは各反復の末尾に先頭へ戻る分岐を持つが，この分岐を if 変換しても
役には立たない．これを除去するということは，ループが取りうる反復回数と同じ
だけ本体を実行するということであり，各コピーには，その反復が実行される
べきだったかどうかを示す述語がそれぞれ必要になる．これらの述語を計算し適用
する命令のコストは，スケジューリングで得られる利得を上回る．さらに，ループ
の分岐は正確に予測され，これにプレディケーションを適用すると，ループの後の
コードを毎反復実行することになる（\cref{tab:l05-which}）．したがってループ
の分岐は残り，命令スケジューリングを 1 回の反復の中に閉じ込める．幅 4 の
インオーダプロセッサで 1 反復 2 サイクルとなる \cref{ex:l10-loop} の
スケジューリング後の本体は，そのプロセッサにとって最良の順序である．本体が
短い場合，ループの独立な命令の大半は異なる反復にあり，そこではコンパイラは
それらを近づけられない．ループ展開がこの制限に対処する．

\section{ループ展開}
\label{sec:l10-unrolling}

本体の短いループでは，実行される命令の大きな割合がループのオーバーヘッドで
あり，本来の処理はループの分岐で隔てられた多数の反復に分散する．
\source{Lesson 10，動画 10--15；H\&P §3.2．}\cref{ex:l10-loop} のループでは，
5つの命令のうち 3つが本来の処理を行い，ADDI と BNE はオーバーヘッドである．

\begin{definition}[ループ展開]
  \label{def:l10-unrolling}
  本体を $k$ 回繰り返して添字を調整し，添字を $k$ ずつ進めることによって，
  新しいループの 1 反復が元のループの $k$ 反復分の処理を行うようにループを
  書き換えることを\term[るーぷてんかい]{ループ展開}[loop unrolling]という．
  この $k$ を\emph{展開係数}という．1 回展開すれば $k=2$，2 回展開すれば
  $k=3$ である．
\end{definition}

$n$ が偶数のとき，$k=2$ で展開するとループは
\begin{lstlisting}[language=C, numbers=none]
for (i = 0; i < n; i += 2) {
    a[i]   = a[i]   + x;
    a[i+1] = a[i+1] + x;
}
\end{lstlisting}
となり，次のように翻訳される．
\begin{lstlisting}[language={[x86masm]Assembler}, numbers=none]
Loop: LW   R3, 0(R1)      ; a[i] <- a[i] + x
      ADD  R3, R3, R2
      SW   R3, 0(R1)
      LW   R3, 4(R1)      ; a[i+1] <- a[i+1] + x
      ADD  R3, R3, R2
      SW   R3, 4(R1)
      ADDI R1, R1, 8      ; 2 要素先へ
      BNE  R1, R5, Loop
\end{lstlisting}
各反復は 8つの命令で 2 要素を処理し，そのうち添字の増分と分岐は 1 回ずつ
しか実行されない．\margin{$n$ が $k$ の倍数でない場合，コンパイラは残りの
要素のためのコード---通常は元のループのコピー---を追加する．}本体の 2つの
コピーは隣接していて間に分岐がないので，2つ目のコピーが中間結果を自分専用
のレジスタに保持するようにすれば，コンパイラは両者をまとめてスケジューリング
できる（\cref{sec:l10-unrolling-cpi}）．

\subsection{命令数と ILP}
\label{sec:l10-unrolling-ilp}

展開はクロック周期には影響しないが，性能の鉄則（\cref{eq:iron-law}）の他の
2つの因子を変える．元のループは 1 要素当たり 5 命令を実行し，展開後の
ループは 2 要素当たり 8 命令，すなわち 1 要素当たり 4 命令を実行する．
オーバーヘッドの命令が実行される回数は半分になり，命令数は
\SI{20}{\percent} 減る．係数~$k$ では，ループは 1 要素当たり $(3k+2)/k$
命令を実行し，これは本来の処理だけの 3 に近づく．\tip{$3 + 2/k$ と読めばよい．
1 要素当たりのオーバーヘッドは $k$ を倍にするごとに半分になるので，利得の大半は
最初の数段階にある．}

ILP も上がる．理想プロセッサ（\cref{eq:l06-cycle}）では，連続する反復の
ADDI が 1 反復 1 サイクルの連鎖をなす．各反復の LW は前の反復の ADDI から
R1 を読み，自分の反復の ADDI と同じサイクルに実行され，その ADD と SW は
続く 2 サイクルに実行される．展開後の本体における R3 への 2 回目の書き込み
は WAW 依存であり，リネーミングがこれを除去する．したがって理想プロセッサ
は 1 サイクルに 1 反復，すなわち $k$ 要素を実行し，長いループの ILP は
1 反復当たりの命令数に近づく．元のループでは 5，展開後のループでは 8 で
ある．これは，反復が互いに添字を通じてしか依存しないことを前提としている．
\texttt{s = s + a[i]} のように累算するループでは，\texttt{s} への加算が
どちらの版でも 1 要素 1 サイクルの連鎖をなす．\tip{このような累算は，$k$ 個の
独立なアキュムレータに分けて展開し，ループの後で合計する．浮動小数点では
\cref{sec:l10-thr}の再グループ化の許可が必要である．}

\subsection{命令当たりサイクル数}
\label{sec:l10-unrolling-cpi}

CPI も下がるかどうかはプロセッサに依存する．\cref{tab:l10-unrolling} は，
\cref{ex:l10-loop} の幅 4 のインオーダプロセッサと理想プロセッサ（どちらも
分岐予測は完全とする）について，ループの 4つの版を比較したものである．
\margin{ループの分岐の予測が外れるのは，実行 1 回につきおよそ 1 回，ループ
を抜けるときである．端数処理のループがあっても高々もう 1 回増えるだけで
ある．}最初の 2つの版は \cref{ex:l10-loop} のものである．展開しただけで
スケジューリングしていない版では，ループは依然として LW--ADD--SW の連鎖を
順に実行する．2つ目の LW が 1つ目の SW と，ADDI が 2つ目の SW と，BNE が
次の反復の LW と一緒に実行されるので，8 命令からなる 1 反復は 5 サイクル
かかる．CPI は 0.6 から $5/8 = 0.625$ に上がるが，1 要素当たりの時間は
命令数とともに 3 サイクルから 2.5 サイクルに減る．理想サイクルの順に
スケジューリングすると，展開後の本体は次のようになる．
\begin{lstlisting}[language={[x86masm]Assembler}, numbers=none]
Loop: LW   R3, 0(R1)
      LW   R4, 4(R1)      ; 2つ目のコピーは R4 を使う
      ADDI R1, R1, 8
      ADD  R3, R3, R2
      ADD  R4, R4, R2
      SW   R3, -8(R1)     ; R1 はすでに進んでいる
      SW   R4, -4(R1)
      BNE  R1, R5, Loop
\end{lstlisting}
\cref{sec:l10-scheduling}の補正が両方とも必要になる．2つ目の LW は，
1つ目のコピーの ADD と SW より前に来るが，これらはまだ R3 を必要とするので，
2つ目のコピーは自分の値を R4 で計算する．また両方のストアが増分の後に来る
ので，その変位を 8 だけ減らす．プロセッサは最初のサイクルに 2つの LW と
ADDI を，2 サイクル目に 2つの ADD を実行し，3 サイクル目を 2つの SW，BNE，
そして次の反復の最初の LW が埋める．それ以降はグループが反復と一致しなく
なり，定常状態では 16 命令，すなわち 2 反復が 5 サイクルかかる．1 要素当たり
1.25 サイクル，CPI は $5/16 = 0.3125$ である．スケジューリングした元の
ループと比べると，1 要素当たりの時間は 2 サイクルから 1.25 サイクルに減り，
高速化率は 1.6 となる．そのうち $5/4 = 1.25$ が命令数に，
$0.4/0.3125 = 1.28$ が CPI に由来する．

\begin{table}[htbp]
  \centering\small
  \caption{$k=2$ で展開する前と後の \cref{ex:l10-loop} のループを，
    コンパイルしたままの版とスケジューリングした版について比較する．
    1 要素当たりの命令数と，分岐予測が完全な 2つのプロセッサにおける
    1 要素当たりのサイクル数および CPI を示す．}
  \label{tab:l10-unrolling}
  \begin{tabular}{@{}lccccc@{}}
    \toprule
    & & \multicolumn{2}{c}{幅 4 のインオーダ} & \multicolumn{2}{c}{理想} \\
    \cmidrule(lr){3-4}\cmidrule(l){5-6}
    版 & 命令数 & サイクル & CPI & サイクル & CPI \\
    \midrule
    元のループ               & 5 & 3.00 & 0.600  & 1.00 & 0.200 \\
    スケジューリング後       & 5 & 2.00 & 0.400  & 1.00 & 0.200 \\
    展開後                   & 4 & 2.50 & 0.625  & 0.50 & 0.125 \\
    展開・スケジューリング後 & 4 & 1.25 & 0.3125 & 0.50 & 0.125 \\
    \bottomrule
  \end{tabular}
\end{table}

このようにインオーダプロセッサでは，展開はそれが可能にするスケジューリング
を通じて CPI を下げる．ループの分岐が 2つの反復の命令を隔てなくなり，
コンパイラが両者を交互に並べられるからである．理想プロセッサでは命令の順序
は問題にならない．そこで CPI が下がるのは ILP が 5 から 8 に上がるから
であり，命令数の減少と合わせて 1 要素当たりの時間は半分になる．大きな
ウィンドウを持つアウトオブオーダプロセッサも，自身の幅までは同じように
利益を得る．1 サイクルに 1 命令を完了するプロセッサは，どの版でも CPI が~1
であり，命令数による 1.25 倍しか得られない．したがってループ展開は，
どのプロセッサでも命令数を下げ，1 サイクルに複数の命令を実行するプロセッサ
では CPI も下げる．

\subsection{展開の欠点}
\label{sec:l10-bloat}

本体は $k$ に比例して大きくなり，5 命令から $3k+2$ 命令になるが，利得は
頭打ちになる．1 要素当たりの命令数は~3 に近づき，実在のプロセッサは ILP が
どれほど高くても IPC を自身の幅より上げられない．$k=4$（14 命令，1 要素
当たり 3.5 命令）から $k=8$（26 命令，1 要素当たり 3.25 命令）にすると，
本体はほぼ 2 倍になるのに，実行される命令は \SI{7}{\percent} しか減らない．
これは第2章の収穫逓減の一例である．

\begin{definition}[コード肥大化]
  \label{def:l10-bloat}
  ループ展開のような変換によってプログラムのサイズが増大することを
  \term[こーどひだいか]{コード肥大化}[code bloat]という．
\end{definition}

コードが大きくなればより多くのメモリを占め，キャッシュの章で見るように，
その命令をフェッチする際のキャッシュミスが増えて，利得を打ち消しうる．
\margin{まず効いてくるのはフロントエンドである．Skylake はデコード済みマイクロ
オペレーションを 1536 個，Zen~2 は 4096 個キャッシュし，それを超えて展開した
ループはデコーダに戻る．Intel 最適化マニュアル．}%
さらに 2つの問題が反復回数に関わる．$n$ が $k$ の倍数でない場合，残りの
反復のために追加のコードが必要になる．また展開は，展開後のループの 1 反復
が始まる時点で残りの反復回数が分かっていることを前提とする．値が 0 である最初の
要素で止まる探索のように，脱出が本体の中で決まるループでは，単純に 1 反復
当たり本体の $k$ 個のコピーを実行することはできない．後ろのコピーが，元の
ループなら実行しない反復を行ってしまう可能性があるからであり，追加の脱出
判定が必要になる．そのためコンパイラは，小さな係数で，しかも本体の短い
ループだけを展開する．

\begin{keypoint}
  ループ展開は，添字の増分とループの分岐を元のループの $k$ 反復につき 1 回
  だけ実行するので，どのプロセッサでも命令数を下げる．さらに，1 サイクルに
  複数の命令を実行するプロセッサでは CPI も下げる．インオーダプロセッサでは
  コンパイラが複数の反復の命令をまとめてスケジューリングできるようになる
  からであり，並列性を自分で見つけるプロセッサでは，反復が独立なループの
  ILP が上がるからである．その代償は，$k$ とともに増大するコードと，反復
  回数が $k$ の倍数でない場合や事前に分からない場合の追加のコードである．
\end{keypoint}

\section{インライン展開}
\label{sec:l10-inlining}

分岐が if-then-else の 2つの経路を隔てるのと同じように，呼び出しは
呼び出し側の命令と呼び出される関数の命令とを隔てる．
\source{Lesson 10，動画 16--18；H\&P 付録 A．}さらに，呼び出しのたびに，
関数本来の処理を何も行わない命令が実行される．呼び出しとリターンそのもの，
呼び出し規約で定められたレジスタへの引数のコピーとそこからの結果のコピー，
そしてしばしばスタックへのレジスタの退避と復帰である．\margin{実際の呼び出し
規約はもっと多くをレジスタで渡す．System~V AMD64 ABI は整数引数の最初の 6つ，
AArch64 は最初の 8つであり，残りはスタックに積まれる．}

\begin{definition}[インライン展開]
  \label{def:l10-inlining}
  呼び出しを，呼び出される関数の本体を呼び出し側のオペランドに合わせた
  コピーで置き換えることを
  \term[いんらいんてんかい]{インライン展開}[function call inlining]という．%
  \caution{C と C++ の \texttt{inline} はリンケージの指定であって命令ではない．
  GCC は \texttt{-O2} 以降，\texttt{inline} と書かれていない関数も展開し，強制
  できるのは \texttt{always\_inline} だけである．}
\end{definition}

\begin{example}
  \label{ex:l10-inlining}
  次の関数を
  \begin{lstlisting}[language=C, numbers=none]
int sqdiff(int x, int y) { int d = x - y; return d * d; }
\end{lstlisting}
  \texttt{c = sqdiff(a, b)} として呼び出す．$a$，$b$，$c$ は R8，R9，R10 に
  ある．呼び出し規約は引数を R4 と R5 で渡し，結果を R2 で返す．
  \begin{lstlisting}[language={[x86masm]Assembler}, numbers=none]
        ADD  R4, R8, R0     ; x <- a
        ADD  R5, R9, R0     ; y <- b
        CALL SqDiff
        ADD  R10, R2, R0    ; c <- 結果
        ...
SqDiff: SUB  R2, R4, R5     ; d <- x - y
        MUL  R2, R2, R2     ; 結果 <- d * d
        RET
\end{lstlisting}
  インライン展開すると，呼び出しは次のようになる．
  \begin{lstlisting}[language={[x86masm]Assembler}, numbers=none]
        SUB  R10, R8, R9    ; c <- a - b
        MUL  R10, R10, R10  ; c <- c * c
\end{lstlisting}
  呼び出しは 1 回につき 7 命令を実行するのに対し，インライン展開後のコード
  は 2 命令である．引数と結果のコピーは，$a$ と $b$ から $c$ への依存の連鎖
  も長くする．理想プロセッサでは，呼び出しは 4 サイクル，インライン展開後の
  コードは 2 サイクルかかる（\cref{fig:l10-inlining}）．
\end{example}

\begin{figure}[htbp]
  \centering
  % Function call inlining: (a) the call with its argument and result copies,
% CALL and RET (shaded: instructions that only implement the call);
% (b) the inlined body.  Circled numbers: cycle in which the ideal processor
% executes each instruction.
\begin{tikzpicture}[x=1cm,y=1cm, font=\footnotesize\sffamily, >={Stealth[length=1.6mm]},
  ins/.style={draw, fill=white, minimum width=2.5cm, minimum height=0.5cm,
    inner sep=2pt, font=\scriptsize\ttfamily},
  ovh/.style={ins, fill=ln-box, draw=ln-rule, text=ln-muted},
  cyc/.style={circle, draw=ln-accent, text=ln-accent, inner sep=0pt,
    minimum size=0.38cm, font=\scriptsize\sffamily},
  ctl/.style={->, thick, dashed, ln-muted},
  lab/.style={font=\scriptsize\sffamily, text=ln-muted}]
  % ---------------- (a) with a call
  \node[align=center] at (1.75,0.95)
    {\iflangja{(a) 呼び出し：7 命令，4 サイクル}{(a) call: 7 instructions, 4 cycles}};
  \node[lab] at (0,0.45) {\iflangja{呼び出し側}{caller}};
  \node[ovh] (c1) at (0,0.0)  {ADD  R4, R8, R0};
  \node[ovh] (c2) at (0,-0.6) {ADD  R5, R9, R0};
  \node[ovh] (c3) at (0,-1.2) {CALL SqDiff};
  \node[ovh] (c4) at (0,-2.4) {ADD  R10, R2, R0};
  \node[lab] at (3.5,-0.15) {SqDiff};
  \node[ins] (f1) at (3.5,-0.6) {SUB  R2, R4, R5};
  \node[ins] (f2) at (3.5,-1.2) {MUL  R2, R2, R2};
  \node[ovh] (f3) at (3.5,-1.8) {RET};
  \foreach \n/\c in {c1/1, c2/1, c3/1, c4/4}
    \node[cyc] at ([xshift=-0.3cm]\n.west) {\c};
  \foreach \n/\c in {f1/2, f2/3, f3/2}
    \node[cyc] at ([xshift=0.3cm]\n.east) {\c};
  \draw[ctl] (c3.east) -- ++(0.35,0) |- (f1.west);
  \draw[ctl] (f3.west) -- ++(-0.35,0) |- (c4.east);
  % ---------------- (b) inlined
  \node[align=center] at (7.3,0.95)
    {\iflangja{(b) インライン展開後：2 命令，2 サイクル}{(b) inlined: 2 instructions, 2 cycles}};
  \node[ins] (i1) at (7.3,-0.6) {SUB  R10, R8, R9};
  \node[ins] (i2) at (7.3,-1.2) {MUL  R10, R10, R10};
  \foreach \n/\c in {i1/1, i2/2}
    \node[cyc] at ([xshift=0.3cm]\n.east) {\c};
\end{tikzpicture}

  \caption{インライン展開の前と後の \cref{ex:l10-inlining}．網掛けの命令は
    呼び出しを実現するだけの命令であり，破線の矢印は制御の移動を，丸で
    囲んだ数字は理想プロセッサが各命令を実行するサイクルを表す．}
  \label{fig:l10-inlining}
\end{figure}

このようにインライン展開は，呼び出し，リターン，引数と結果のコピーがなく
なるので命令数を下げ，本体が呼び出し側のオペランドを直接読み，結果を
呼び出し側が必要とする場所に直接届けるので依存の連鎖を短くする．さらに，
本体は呼び出し側のコードの一部になる．コンパイラはその命令を呼び出し側の
命令とまとめてスケジューリングでき，かつての呼び出しをまたいで木の高さ削減
を適用でき，プロセッサは呼び出しとリターンを予測する必要がなくなる
（第4章）．\sidenote{インライン展開は何よりも，他の最適化を\emph{可能にする}
変換である．本体が呼び出し側に入ると，引数はそこではしばしば定数として分かって
いるので，定数伝播やデッドコード除去，そしてこの章の技法が，かつての境界をまたい
で適用できる．そのためコンパイラは早い段階でインライン展開し，その後で最適化する．}

\subsection{インライン展開によるコード肥大化}
\label{sec:l10-inlining-bloat}

インライン展開は本体のコピーを呼び出し箇所ごとに別々に置くので，多数の箇所
から呼ばれる大きな関数ではコードが相当に増える．

\begin{example}
  \label{ex:l10-inlining-size}
  本体が 20 命令とリターンからなる関数が 10 箇所から呼ばれ，各箇所が
  呼び出しに 4 命令（引数のコピー 2つ，呼び出し，結果のコピー 1つ）を
  費やすとする．インライン展開しなければ，コードは
  $10 \times 4 + 21 = 61$ 命令を占め，呼び出し 1 回につき $4 + 21 = 25$
  命令が実行される．すべての箇所でインライン展開すると，関数と呼び出しの列
  は消えるが，各箇所が本体の 20 命令を受け取る．合計 200 命令，コードは
  3 倍以上になるのに，節約できるのは呼び出し 1 回当たり 25 命令のうち
  5 命令（\SI{20}{\percent}）である．\margin{1 次命令キャッシュは小さいままである．Skylake も Zen~3 も
  32~KB であり，インライン展開されたコピーはそこを奪い合う．}本体が 2 命令の \cref{ex:l10-inlining}
  の関数なら，同じ 10 箇所でも，インライン展開しない場合は
  $10 \times 4 + 3 = 43$ 命令，展開した場合は $10 \times 2 = 20$ 命令しか
  必要としない．コードは小さくなり，呼び出し 1 回につき実行される命令は
  7 から 2 になる．
\end{example}

ループ展開と同様に，コードが大きくなるとキャッシュミスが増え，利得を
打ち消しうる．そのためコンパイラがインライン展開するのは，本体に比べて
呼び出しのオーバーヘッドが大きい小さな関数と，呼び出し箇所が 1つだけで
そのコピーが元の関数を置き換える関数である．多数の箇所から呼ばれる大きな
関数は，呼び出しのまま残る．

\begin{keypoint}
  インライン展開は，呼び出し，リターン，引数と結果の受け渡しを除去して
  命令数を下げ，依存の連鎖を短くし，さらにコンパイラが本体を呼び出し側と
  まとめて最適化できるようにする．小さな関数では効果があるが，多数の箇所
  から呼ばれる大きな関数では，コード肥大化が利得を上回る．
\end{keypoint}

\section{その他のコンパイラ技法}
\label{sec:l10-other}

この章の技法は，IPC を高めるコンパイラ変換という，より大きな体系の例で
ある．\source{Lesson 10，動画 19--20；\cite{lam1988}；\cite{fisher1981}．}%
命令スケジューリングを，それを閉じ込めている境界の向こうまで押し広げる技法
がさらに 2つあるが，ここでは概略を述べるにとどめる．
\term[そふとうぇあぱいぷらいにんぐ]{ソフトウェアパイプライニング}%
[software pipelining]では，コンパイラはループをパイプラインのように扱う．
反復が重なるように連続する反復の命令をスケジューリングし，各サイクルが複数
の反復の独立な命令を組み合わせつつ，結果は元のループのものと同じに保たれる
\cite{lam1988}．もう 1つの
\term[とれーすすけじゅーりんぐ]{トレーススケジューリング}%
[trace scheduling]は，if 変換をさらに推し進める．コンパイラはプロファイルを
用いて，複数の分岐を通る最も頻繁に実行される経路を選び，この経路を 1つの
コードブロックとしてスケジューリングし，実行がその経路から外れたときに結果を
訂正する補償コードへ分岐する判定を追加する \cite{fisher1981}．コンパイラに
よるスケジューリングを突き詰めると，操作を並列なステップにグループ化すること
を完全にコンパイラに委ねるプロセッサ，すなわち第11章の VLIW プロセッサに
行き着く．\sidenote{その VLIW はソフトウェアパイプライニングをハードウェアで
支援していた．IA-64 の\emph{回転}レジスタは反復ごとにレジスタを付け替え，回転
述語がパイプラインの各段を有効・無効にし，専用のループ分岐が立ち上げと立ち下げを
行う．そのため，展開なら複数のコピーを要する処理を本体 1つで行える．}

\begin{keypoint}[章のまとめ]
  IPC は，プログラムの ILP と，プロセッサがウィンドウの中で見つけられる
  独立な命令とによって制限され，コンパイラはそのどちらにも影響する．木の
  高さ削減は式を再グループ化してその命令がより短い連鎖をなすようにし，
  命令数を変えずに ILP を上げる．命令スケジューリングは独立な命令を隣り
  合わせにし，命令が互いを追い越す箇所では宛先レジスタとアドレスのオフ
  セットを変更する．これはインオーダプロセッサやウィンドウの小さい
  プロセッサで効く．分岐はこれを閉じ込めるが，if 変換は小さな if-then-else
  から分岐を除去する．ループ展開は，添字の増分とループの分岐を $k$ 反復に
  つき 1 回だけ実行して命令数を下げる．コンパイラが複数の反復の命令をまとめ
  てスケジューリングできるようになることと，反復が独立なループの ILP が
  上がることから，CPI も下げる．インライン展開は，呼び出し，リターン，引数
  のコピーを除去する．どちらもコード肥大化を招くので，コンパイラは短い
  ループと小さな関数に限ってこれらを用いる．ソフトウェアパイプライニングと
  トレーススケジューリングは，スケジューリングをループの反復をまたいで，
  また分岐をまたいで拡張する．
\end{keypoint}

\end{document}
